\documentclass{article}
\usepackage[utf8]{inputenc}
\usepackage{amsmath}
\usepackage{geometry}
\geometry{margin=2cm}
\begin{document}

\section*{Análise Completa da Questão sobre Mediana da Umidade Relativa do Ar}

\subsection*{Enunciado}
A umidade relativa do ar é um dos indicadores utilizados na meteorologia para fazer previsões sobre o clima. O quadro apresenta as médias mensais, em porcentagem, da umidade relativa do ar em um período de seis meses consecutivos em uma cidade. Nesta cidade, a mediana desses dados, em porcentagem, da umidade relativa do ar no período considerado foi:

\begin{center}
\begin{tabular}{|c|c|c|c|c|c|c|}
\hline
\cellcolor{blue!20}\textbf{Meses} & Maio & Jun. & Jul. & Ago. & Set. & Out. \\
\hline
\cellcolor{blue!20}\textbf{Média mensal} \\ \cellcolor{blue!20} da umidade relativa do ar (\%)} & 66 & 64 & 54 & 46 & 60 & 64 \\
\hline
\end{tabular}
\end{center}

\subsection*{Solução Técnica}
\begin{itemize}
\item Dados: 66, 64, 54, 46, 60, 64.
\item Ordenação crescente: 46, 54, 60, 64, 64, 66.
\item Número de elementos (n) = 6, par.
\item Cálculo da mediana:
\[
\text{Mediana} = \frac{\text{elemento}_3 + \text{elemento}_4}{2} = \frac{60 + 64}{2} = 62.
\]
\item Portanto, a mediana da umidade relativa do ar no período é 62\%.
\end{itemize}

\subsection*{Formulação Matemática Utilizada}
\begin{equation*}
\text{Mediana em dados pares} = \frac{ \text{Valor no índice} \frac{n}{2} + \text{Valor no índice} \left(\frac{n}{2} + 1\right)}{2}.
\end{equation*}

\subsection*{Explicação Didática}
A mediana representa o valor central que separa os dados ordenados em duas partes iguais. Quando a quantidade de dados é par, ela não coincide com um dado específico, mas é a média dos dois valores centrais.

\subsubsection*{Passos simplificados}
\begin{enumerate}
\item Anote os dados conforme fornecido.
\item Ordene os dados em ordem crescente.
\item Identifique os dois elementos centrais.
\item Calcule a média desses valores centrais.
\item A média obtida é a mediana.
\end{enumerate}

\subsubsection*{Erros comuns}
\begin{itemize}
\item Não ordenar os dados corretamente.
\item Usar apenas um elemento central para calcular a mediana em conjuntos com elementos pares.
\item Escolher incorretamente os valores centrais.
\item Fazer arredondamento inadequado.
\end{itemize}

\subsection*{Dicas para o Professor}
\begin{itemize}
\item Reforce a necessidade de ordenar os dados antes de qualquer cálculo.
\item Explique a diferença no cálculo da mediana para conjuntos com números pares e ímpares de elementos.
\item Use exemplos práticos que envolvam dados estatísticos reais para facilitar a compreensão.
\item Estimule a utilização de calculadoras para precisão.
\item Contextualize a mediana com outras medidas estatísticas.
\end{itemize}

\subsection*{Análise dos Distratores}
\begin{description}
\item[56] Possível erro de ordenação ou escolha incorreta dos elementos centrais; parece apropriado por estar próximo à média, mas incorreto.
\item[58] Erro no cálculo da média entre valores centrais; valor plausível, porém incorreto.
\item[59] Erro de arredondamento ou seleção incorreta; valor intermediário que não corresponde à mediana certa.
\item[62] Resposta correta, baseada no cálculo exato da mediana.
\end{description}

\subsection*{Perfil da Questão}
\begin{itemize}
\item Habilidade principal: Interpretação e análise de dados estatísticos.
\item Habilidades secundárias: Organização e ordenação de dados; cálculo da mediana.
\item Grau de dificuldade: Médio.
\item Necessita interpretação visual e cálculo numérico.
\end{itemize}

\subsection*{Revisão de Consistência}
A análise da questão está consistente e correta, exceto pela correção necessária no valor da resposta oficial, que deve ser 62 em vez de 60 para refletir os cálculos matemáticos.

\vspace{1em}

\textbf{Fim da análise.}

\end{document}